%% annex/p02/first-order-completeness-remainder.tex --- retained proofs
\noindent\textit{Reader note.} The exact inherited proofs through source line 6048
replace historical annex items 2.2.17--2.2.19.  The following exact historical
lines 303--338 support only the separately attributed elementary-diagram
additions and do not replace inherited proofs.

\begin{annexitem}{annex:p2.2.20}{thm:p02-upward-ls-elementary-extension}{Upward L\"owenheim--Skolem via elementary diagrams}
\begin{proof}
Expand \(\Lang(\mathcal M)\) further by constants
\(d_\alpha\) for \(\alpha<\kappa\), and adjoin to
\(\operatorname{ElDiag}(\mathcal M)\) all
\(d_\alpha\ne d_\beta\) for \(\alpha<\beta<\kappa\) and all
\(d_\alpha\ne c_a\) for \(a\in\dom\mathcal M\).  Every finite fragment is
satisfied by the parameter expansion of \(\mathcal M\): because its domain
is infinite, interpret the finitely many occurring \(d_\alpha\) as distinct
elements avoiding the finitely many named \(a\)'s.  Compactness supplies a
model \(\mathcal P\) of the whole theory.

Let \(A\) comprise the interpretations of every \(c_a\) and
\(d_\alpha\).  It has cardinality \(\kappa\).  Downward
L\"owenheim--Skolem in the full expanded signature, whose cardinality is
\(\kappa\), gives
\(\mathcal Q\preccurlyeq\mathcal P\) containing \(A\) and of size at most
\(\kappa\), hence exactly \(\kappa\).  Let \(\mathcal N\) be its
\(\Lang\)-reduct and define \(j(a)=c_a^{\mathcal Q}\).  Distinctness
sentences make \(j\) injective.  For each parameter sentence, it or its
negation belongs to the elementary diagram according to its truth in
\(\mathcal M\), so \(j\) preserves and reflects every \(\Lang\)-formula.
Thus \(j\) is elementary.  The interpretation of
\(d_0\) is outside its range, making the range proper.
\end{proof}
\end{annexitem}

\begin{annexitem}{annex:p2.2.21}{thm:p02-nonstandard-arithmetic-model}{A countable nonstandard elementary model of arithmetic}
\begin{proof}
Apply the preceding theorem to the countable infinite structure
\(\mathcal N_{\mathrm{ar}}\) with \(\kappa=\aleph_0\).  Replacing the target
by an isomorphic copy over the elementary image makes that image an actual
proper substructure.  The resulting countable elementary extension contains
an element outside every embedded standard numeral.
\end{proof}
\end{annexitem}
